现在，你坐在哨兵的脑袋里面，你的目标是悄悄摸到敌方基地后面偷家。
但是没认真看规则的你忘记了地图长什么样，你打开手机，发现好心的wzx部长已经帮你存了一份\textbf{标准赛场地图}。
于是你打开\textbf{缺德地图}，把赛场地图传上去。你又打开手机GPS，把自己的\textbf{定位}也告诉缺德地图，
最后，你输入目的地，缺德地图就帮你规划好了一条用时最少的路。
于是你就跟着这条路径走，同时你并没有关掉GPS，这样才知道自己有没有走错。
在路上，你在\textbf{视野}里看到了许多其他机器人，你以你秋名山车神的\textbf{车技}，在他们之间穿梭，避开了所有车。
突然！What can I say, man!无人机坠机坠在了你前进的道路上。
你对车身大小了如指掌。只看一眼，你计算出哨兵过不去，路被堵了。
鉴于无人机死状惨烈，这无人机短时间是飞不起来，你走不了这条路了。
不过还好，每隔一段时间，缺德地图会再算一次路径，你可能会听到“正在为您重新规划路径”，他帮你找到了另一条路。
虽然“缺德地图持续为您服务”，但是有时候你可能会开到一些“信号不好”，无法获取缺德地图服务的地方。
此时，你要能够在不依靠缺德地图的情况下\textbf{自己找到}信号比较好能重新连接上缺德地图服务的地方。
最终你成功到达了目的地。

从上面可以看出，导航需要：
\textbf{[1]地图（全局代价地图）}
\textbf{[2]导航系统（全局规划器）}
\textbf{[3]定位系统（多传感器融合）}
\textbf{[4]附近的视野（局部代价地图）}
\textbf{[5]车手（局部规划器）}
\textbf{[6]应急方案（恢复行为树）}